\scanpage{36}
\[
(18)\qquad \varphi:\ p^*(E_{Y,X})\longrightarrow i^*(E_{X^*}).
\]

En termes du diagramme de groupes
(10) Quand les \inserted{trois} groupoïdes composants de
(8) sont connexes, i.e. $Y$, $X^*$ et $\mathcal V^*_{Y,X}$
connexes), on peut aussi interpréter
ces systèmes, \deleted{\unsure{comme des systèmes}}\inserted{en considérant $E_{X^*}$ comme}
un ens. sur lequel $\Pi_1(X^*)$ opère, $E_{Y,X}$
comme un ens. sur lequel $\Pi_1(Y)$ opère,
et $p^*(E_{Y,X})$ et $i^*(E_{X^*})$ dans (18) comme
étant ces $=$s ensembles, \deleted{\unsure{munis}}\inserted{considérés} comme
$\Pi_1(\mathcal V^*_{Y,X})$-ensembles \deleted{qui} par restriction
des opérateurs.

\margin{NB Dans le cas où on a $=$ un
diagramme (14), cela veut que $\Pi_1(Y)$
apparaît comme un quotient de $D=\Pi_1(\mathcal V^*_{Y,X})$, et
donc $E_{Y,X}$ peut s'interpréter d'emblée comme
un $\deleted{\unsure{\Pi_1(\mathcal V^*_{Y,X})}}D$-ens., sur lequel le ss-gr.
$I$ opère trivialement, on
interprète $\varphi$ de (18) comme un
hom. de $\Pi_1(Y)$-ens.
\[
(18\ \mathrm{bis})\qquad E_{Y,X}\longrightarrow
\underbrace{E^I_{X^*}}_{I\text{-invariants}}
\]}

Un cas \inserted{typique,} particulièrement important
pour \unsure{la|les} \deleted{applic. à} Géométrie algébrique,
est celui où $X$ est une surface
top. et où $Y$ est réduit à un pt $a$,
donc $\mathcal V_{Y,a}=\mathbb D_a$, \uline{disque} \inserted{(infinitésimal)} autour de $a$,
et $\mathcal V^*_{Y,a}=\mathbb D^*_0$, disque épointé (inf.) autour
de $a$, de sorte qu'on trouve le diagramme
% Pale bottom-left foliation, in another hand; not the author's own numbering.
\margin{35}
